
tenu de la structure de \(K^{\bullet}(\underline{\underline{D}}(\underline{E}))\) sur \(K^{\bullet}(S)\), déterminée par 4.7, et du fait que les \(\xi_i^{(j)}\) sont les classes des faisceaux intervenant dans le scindage de l'image inverse de \(\underline{E}\) sur \(\underline{\underline{D}}(\underline{E})\). Cela achève la démonstration de la première forme de 4.6.

Pour en déduire la deuxième forme, on remarque que les éléments \(F_i\) de la \(K^{\bullet}(S)\)-\(\lambda\)-algèbre \(K^{\bullet}(X)\) vérifient les relations i) et ii). Si \(A\) est la \(K^{\bullet}(S)\)-\(\lambda\)-algèbre engendrée par des éléments \(F_i'\) satisfaisant i) et ii), il existe donc un \(K^{\bullet}(S)\)-\(\lambda\)-homomorphisme de \(A\) dans \(K^{\bullet}(X)\) envoyant \(F_i'\) sur \(F_i\) (V 4.7). D'après 4.6, c'est homomorphisme est surjectif. Comme le \(\lambda\)-idéal engendré par les relations i) et ii) dans \(K^{\bullet}(S)[F_1,\ldots,F_k]\) contient l'idéal usuel engendré par les relations de 4.6, cet homomorphisme est aussi injectif, donc est un isomorphisme.

\textbf{Corollaire 4.8.} \uline{Avec les hypothèses et les notations de} 4.5, \uline{posons}~:
\[
  c_i^{(j)}=\gamma^i(F_j-p_j),
  \qquad 1\le i\le p_j;
\]
\uline{alors la} \(K^{\bullet}(S)\)-\uline{algèbre} \(K^{\bullet}(X)\) \uline{est engendrée par les} \(c_i^{(j)}\); \uline{l'idéal des relations entre les} \(c_i^{(j)}\) \uline{est engendré par les relations déduites de}~:
\[
  \gamma_t(E-n)=\prod_{j=1}^{k}\gamma_t(F_j-p_j).
\]

Ce n'est qu'un autre énoncé de 4.6. On peut en effet exprimer les \(c_i^{(j)}\) en fonction des \(\lambda_i^{(j)}\), et réciproquement. De plus, la relation~: \(\lambda_t(E)=\prod\lambda_t(F_j)\) équivaut à la relation~: \(\gamma_t(E-n)=\prod\gamma_t(F_j-p_j)\).

4.9.
Soit \(S\) un schéma; notons \(\underline{\underline{G}}_S^{n,r}\) la grassmanienne de type \((n,r-n)\) de \((\underline{O}_S)^r\), avec \(r\ge n\). Pour \(r\) variable, les \(\underline{\underline{G}}_S^{n,r}\) forment un système inductif, le morphisme \(\underline{\underline{G}}_S^{n,r}\longrightarrow \underline{\underline{G}}_S^{n,r+1}\) étant défini en considérant \((\underline{O}_{\underline{\underline{G}}_S^{n,r}})^r\) comme
